
\textbf{Proposition 3.2.} \uline{Soit} $f:X\to Y$ \uline{un morphisme projectif}, \uline{d'intersection complète}, \uline{et dont la dimension relative virtuelle} $d$ \uline{est constante}. \uline{On suppose que} $Y$ \uline{est quasi-compact}, \uline{et admet un faisceau ample}. \uline{On a alors pour tout} $k\in\mathbb Z$:
\[
f_*\bigl(\Fil^kK^{\bullet}(X)_{\mathbb Q}\bigr)
\subset
\Fil^{k-d}K^{\bullet}(Y)_{\mathbb Q}.
\]

On obtient donc en particulier un homomorphisme additif :
\[
f_{\mathrm{gr}}:\Gr K^{\bullet}(X)_{\mathbb Q}\longrightarrow
\Gr K^{\bullet}(Y)_{\mathbb Q}
\]
envoyant $\Gr^kK^{\bullet}(X)_{\mathbb Q}$ dans $\Gr^{k-d}K^{\bullet}(Y)_{\mathbb Q}$.

\textbf{3.3.} Rappelons que nous avons défini dans (V 6.1) un foncteur $A\longmapsto \operatorname{Ch}(A)$ sur la catégorie des $K$-algèbres graduées, à valeurs dans la catégorie des $K$-$\lambda$-algèbres augmentées, $K$ étant un anneau binômial (V 2.7). Si $A$ est une $K$-algèbre graduée, on a :
\[
\operatorname{Ch}(A)=K\times(1+\widehat A^+).
\]

Pour toute $K$-$\lambda$-algèbre augmentée $\Lambda$, nous avons défini un $K$-$\lambda$-homomorphisme, nommé «~classe de Chern complétée~» (V 6.7):
\[
\widetilde c=(\varepsilon,c):\Lambda\longrightarrow\operatorname{Ch}(\Gr_{\lambda}\Lambda).
\]

Enfin, nous avons défini un homomorphisme d'anneaux, fonctoriel en $A$, et nommé caractère de Chern (V 6.4):
\[
\ch:\operatorname{Ch}(A)\longrightarrow K\oplus\widehat A^+_{\mathbb Q},
\]
et un endomorphisme de groupes, nommé caractère de Todd (V 1.16):
\[
\Todd:1+\widehat A^+\longrightarrow1+\widehat A^+.
\]
